\scanpage{61}
puisque 2 \deleted{\unsure{(donc 4) est}} (donc 4) est $\underline O_S$-régulier. Donc $\Delta(b,c)\simeq\Delta(b',c')$, et l'isom est unique, cqfd.

b) \deleted{L'argument précédent montre plus généralement, sans hyp. $S$ affine, que la catégorie des $\Delta(b,c)$ correspondant à $(L,D,b_0)$ donnés, est \inserted{$\simeq$ cat.} ponctuelle. Or si la condition $D_1=$ Considérons} $H^1(S,L)\unsure{-|\to}H^1(S,L_0)$ injectif implique que les $\Delta$ correspondants à $(L,D,b_0)$ correspondent à une enveloppe affine qui est un $\underline L$-torseur trivial, i.e. admet un scindage. Ces $\Delta$ sont donc isom. à des $\Delta(b,c)$, \deleted{Or $\unsure{\Delta(b,}$} $b\in\Gamma(S,L)$, $c\in\Gamma(S,L^{\otimes2})$. \deleted{Donc} Pour que $\Delta(b,c)$ donne $b_0,D$, il f et s. que
\[b_{S_0}=b_0,\qquad b^2-4c=D.\]

Donc $(L,D,c_0)$ est dans l'image essentielle ssi $\forall\,b,c$ satisfaisant les conditions précédentes. Ces conditions impliquent évidemment que $D_1=b_1^2$, où $b_1\in\Gamma(S_1,L_1)$ est tel que $b_1^{\deleted{\unsure{*|0}}}|S_0=b_0$.

Inversement, s'il existe un tel $b_1$, \deleted{\unsure{il|à}} les $b$ qui le relèvent forment un torseur sous $L$, \deleted{provient d'une section de $L$, grâce à la s. e.} dont la restriction à $S-S_0$ est trivial, donc \deleted{cette décl. cohomologique associée à $0\to\unsure{L}\xrightarrow{4}L\to L_1\to0$} trivial par l'hypothèse \inserted{sur $S,S_0$}, donc ce torseur admet \deleted{$H^0(S,L)\to H^0(S_1,L_1)\to H^1(S,L)\xrightarrow{4}H^1(S,L)$} une section $b$. Comme $b^2-D\in4\underline L^{\otimes2}$ et que \deleted{Or $H^1(S,L)_2=H^1(S\setminus S_0,L)$} \deleted{\unsure{(2}}4) est $\mathcal O_S$-régulier donc $L^{\otimes2}$-régulier, on peut écrire $b^2-D=\deleted{\unsure{-}}4c$, i.e. $b^2-4c=D$, cqfd

c) La condition d'existence de $\Delta$, \inserted{correspondant à $(L,D,b_0)$,} à cause de a) doit être locale, \deleted{\unsure{il faut}} comme la condition $D_1=\gamma(b_0)$.
% archive sheet number: 60
